\oldpage[0\textsubscript{III}]{67}
les opérateurs de dérivation dans $K^\bullet$ étant limites projectives
des opérateurs correspondants dans les $K_\alpha^\bullet$.

Considérons les suites exactes
\[
(*_n)\qquad
0\longrightarrow \mathrm B^n(K_\alpha^\bullet)
\longrightarrow \mathrm Z^n(K_\alpha^\bullet)
\longrightarrow H^n(K_\alpha^\bullet)\longrightarrow0
\]
\[
(**_n)\qquad
0\longrightarrow \mathrm Z^{n-1}(K_\alpha^\bullet)
\longrightarrow K_\alpha^{n-1}
\longrightarrow \mathrm B^n(K_\alpha^\bullet)\longrightarrow0
\]

L'hypothèse et la prop.
(\hyperref[0.13.2.1-fr]{13.2.1}, (i)) montrent que le système projectif
$(\mathrm B^n(K_\alpha^\bullet))_{\alpha\in I}$ vérifie (ML) pour tout
$n$~; il résulte donc de (\hyperref[0.13.2.2-fr]{13.2.2}) que la suite
\[
(***_n)\qquad
0\longrightarrow\varprojlim_\alpha\mathrm B^n(K_\alpha^\bullet)
\longrightarrow\varprojlim_\alpha\mathrm Z^n(K_\alpha^\bullet)
\longrightarrow\varprojlim_\alpha H^n(K_\alpha^\bullet)
\longrightarrow0
\]
est exacte. Or il est clair que
$\varprojlim\mathrm B^n(K_\alpha^\bullet)$ s'identifie à un
sous-groupe de $K^{n+1}$ contenant $\mathrm B^n(K^\bullet)$ et que
$\varprojlim\mathrm Z^n(K_\alpha^\bullet)$ s'identifie à un
sous-groupe de $\mathrm Z^n(K^\bullet)$~; par suite, $h_n$ est surjective.
Si maintenant, on suppose en outre que le système projectif
$(H^{n-1}(K_\alpha^\bullet))_{\alpha\in I}$ vérifie (ML), les suites
exactes $(*_{n-1})$ et la prop.
(\hyperref[0.13.2.1-fr]{13.2.1}, (ii)) montrent que le système projectif
$(\mathrm Z^{n-1}(K_\alpha^\bullet))_{\alpha\in I}$ vérifie (ML)~; mais
alors, (\hyperref[0.13.2.2-fr]{13.2.2}) appliquée aux suites exactes
$(**_n)$ montre que la suite
\[
0\longrightarrow
\varprojlim_\alpha\mathrm Z^{n-1}(K_\alpha^\bullet)
\longrightarrow K^{n-1}\xrightarrow{u}
\varprojlim_\alpha\mathrm B^n(K_\alpha^\bullet)
\longrightarrow0
\]
est exacte~; comme
$\varprojlim\mathrm B^n(K_\alpha^\bullet)
\supset\mathrm B^n(K^\bullet)$, et que la composée de l'injection
$\varprojlim\mathrm B^n(K_\alpha^\bullet)\to K^n$ et de $u$ est
l'opérateur de dérivation $K^{n-1}\to K^n$, le fait que $u$ soit surjectif
entraîne
$\varprojlim\mathrm B^n(K_\alpha^\bullet)
=\mathrm B^n(K^\bullet)$, donc $h_n$ est injective. C.Q.F.D.
\end{proof}

\begin{env}[13.2.4]
\label{0.13.2.4-fr}
\emph{Remarques} (13.2.4). ---
\begin{enumerate}
\item[(i)]
Le raisonnement de (\hyperref[0.13.2.2-fr]{13.2.2}) (cf. Bourbaki,
\emph{loc. cit.}) montre que la conclusion de cette proposition reste
valable lorsqu'on suppose seulement que les $A_\alpha$ peuvent être munis
de structures d'espaces \emph{métrisables complets}, dans lesquels les
translations sont des homéomorphismes, que les applications
$f_{\alpha\beta}:A_\beta\to A_\alpha$ définissant le système projectif
$(A_\alpha)$ sont \emph{uniformément continues} pour les distances
considérées, et enfin que le système $(A_\alpha)$ vérifie la condition
\begin{itemize}
\item[(ML$'$)]
\emph{Pour tout indice $\alpha$, il existe $\beta\geq\alpha$ tel que pour
tout $\gamma\geq\beta$, $f_{\alpha\gamma}(A_\gamma)$ est dense dans
$f_{\alpha\beta}(A_\beta)$.}
\end{itemize}

Cela permet d'apporter un complément analogue à
(\hyperref[0.13.2.3-fr]{13.2.3})~: supposons que
$K_\alpha^n=0$ pour $n<0$ et pour tout $\alpha$~; supposons de plus que
$(K_\alpha^n)_{\alpha\in I}$ vérifie (ML) pour $n\geq0$ et que les
$A_\alpha=H^0(K_\alpha^\bullet)$ puissent être munis de structures
d'espaces métriques vérifiant les propriétés ci-dessus. Alors les
conclusions de (\hyperref[0.13.2.3-fr]{13.2.3}) sont inchangées pour
$n\geq2$, et en outre $h_1$ est \emph{bijectif}, car le raisonnement de
(\hyperref[0.13.2.2-fr]{13.2.2}) montre encore que
$(\mathrm B^1(K_\alpha^\bullet))_{\alpha\in I}$ vérifie (ML), que la suite
$(***_1)$ est exacte, et enfin, en vertu de ce qui précède, que
$\varprojlim\mathrm B^1(K_\alpha^\bullet)=\mathrm B^1(K^\bullet)$. On a
ainsi établi entre autres les assertions de (T, 3.10.2).

\item[(ii)]
Il est possible d'introduire les foncteurs dérivés à droite
$\varinjlim^{(n)}$ du foncteur $\varinjlim$, et d'obtenir des énoncés
plus complets que les précédents [28].
\end{enumerate}
\end{env}

